
[[1]]
Exercise 1: Infinitesimal Lorentz transformation ( $3+3$ points)
An infinitesimal Lorentz transformation and its inverse can be written as
$$
x^{\prime \alpha}=\left(g^{\alpha \beta}+\epsilon^{\alpha \beta}\right) x_\beta, \quad x^\alpha=\left(g^{\alpha \beta}+\epsilon^{\prime \alpha \beta}\right) x_\beta^{\prime}
$$
with the Minkowski metric $\left(g^{\alpha \beta}\right)=\operatorname{diag}(1,-1,-1,-1)$ and $\epsilon^{\alpha \beta}$ and $\epsilon^{\alpha \beta}$ are infinitesimal.

[[1.1]]
i) Show that $\epsilon^{\prime \alpha \beta}=-\epsilon^{\alpha \beta}$ by making use of the definition of an inverse transformation.

[[1.2]]

ii) Show from the preservation of the norm that the infinitesimal shift $\epsilon^{\alpha \beta}$ is antisymmetric $\epsilon^{\alpha \beta}=-\epsilon^{\beta \alpha}$.


[[2.1]]
Exercise 2: Collider experiments at the LEP (4 points)
In the LEP storage ring at CERN, head-on collisions between (equally accelerated) electrons and positrons were produced, such that the total energy in the center of mass was equal to that of the Z boson ( $m_z=91 \mathrm{GeV}$ ). What is the velocity of each particle before the collision? If an electron is accelerated toward a positron at rest, what velocity does it need in order to reach the same center-of-mass total energy?
Hint: The total momentum of two particles $a, b$ in the center of mass frame is $\mathbf{p}_a+\mathbf{p}_b=0$.

[[3.1]]
Exercise 3: Continuity equation in relativistic quantum mechanics ( $3+2$ points) In the lecture you derived the Klein-Gordon equation of relativistic quantum mechanics
$$
-\frac{\partial^2}{\partial t^2} \phi(x)=\left[-\Delta+m^2\right] \phi(x)
$$

[[3.1]]
i) Show that the current
$$
j^\mu(x)=\mathrm{i}\left(\phi^*\left(\partial^\mu \phi\right)-\left(\partial^\mu \phi^*\right) \phi\right)
$$
satisfies a continuity equation
$$
\partial^\mu j_\mu=0
$$

[[3.2]]
ii) Since the Klein-Gordon equation is a relativistic wave equation, it is solved by plane waves. Calculate the zero-component of the current $j^0=\rho$ for a plane wave solution of the form $\phi \sim \exp (\mathrm{i} p x)$ and interpret your result.


[[4]]
Exercise 4: Lorentz invariant integration measure (5 points)
Given the relativistic invariance of $d^4 k$ show that the integration measure
$$
\frac{d^3 k}{(2 \pi)^3 2 E(\mathbf{k})}
$$
is Lorentz invariant, provided that $E(\mathbf{k})=k_0=\sqrt{m^2+\mathbf{k}^2}$. 

Use this to argue that $2 E(\mathbf{k}) \delta^{(3)}\left(\mathbf{k}-\mathbf{k}^{\prime}\right)$ is a Lorentz invariant distribution.
Hint: Start from the Lorentz-invariant expression $\frac{d^4 k}{(2 \pi)^3} \delta\left(k^2-m^2\right) \theta\left(k_0\right)$ and use $\delta\left(x^2-x_0^2\right)=\frac{1}{2\left|x_0\right|}\left(\delta\left(x+x_0\right)+\delta\left(x-x_0\right)\right)$.

[[4.1]]
Given the relativistic invariance of $d^4 k$ show that the integration measure
$$
\frac{d^3 k}{(2 \pi)^3 2 E(\mathbf{k})}
$$
is Lorentz invariant, provided that $E(\mathbf{k})=k_0=\sqrt{m^2+\mathbf{k}^2}$. 

[[4.2]]
Use this to argue that $2 E(\mathbf{k}) \delta^{(3)}\left(\mathbf{k}-\mathbf{k}^{\prime}\right)$ is a Lorentz invariant distribution.
Hint: Start from the Lorentz-invariant expression $\frac{d^4 k}{(2 \pi)^3} \delta\left(k^2-m^2\right) \theta\left(k_0\right)$ and use $\delta\left(x^2-x_0^2\right)=\frac{1}{2\left|x_0\right|}\left(\delta\left(x+x_0\right)+\delta\left(x-x_0\right)\right)$.
